\documentclass[10pt,journal,compsoc,twocolumn]{IEEEtran}
\usepackage[utf8]{inputenc}
\usepackage{amsmath,amssymb,amsfonts}
\usepackage{graphicx}
\usepackage{booktabs}
\usepackage{cite}
\usepackage{hyperref}
\usepackage{microtype}

\begin{document}

\title{Task A3: Dry Bean Variety Recognition via Morphological Modeling and Dimensionality Reduction}

\author{
    \IEEEauthorblockN{Muhammad Abdullah}\\
    \IEEEauthorblockA{Department of Intelligent Systems, The University of Lahore, Pakistan\\
    Roll No: 70178208 \quad|\quad Section: BSRAI-5C\\
    GitHub: \url{https://github.com/MuhammadaAbdullah70178208/AI101-ML-Assignment-70178208}\\
    Instructor: Dr. Muhammad Awais, AI Architect}
}

\maketitle

\begin{abstract}
Automated optical grain sorting requires precise multi-class classification capable of distinguishing morphologically similar agricultural varieties under class imbalance. This study addresses Task A3 of the AI-101 course curriculum, applying the inductive learning formulation and ML lifecycle from Modules~01 and~02 to the UCI Dry Bean benchmark ($N = 3,500$ grains, 7 varieties, 16 morphological attributes). We evaluate Multinomial Logistic Regression, Balanced Logistic Regression, Support Vector Machines (SVM RBF), Random Forest, and Multi-Layer Perceptrons (MLP). The SVM RBF model achieves 85.52\% test accuracy and 0.8616 Macro-F1, while balanced class weighting improves minority sensitivity to 0.8652 Macro-F1. Furthermore, Principal Component Analysis (PCA) reveals that an 8-dimensional subspace preserves 99.15\% of discriminative performance, validating efficient lower-dimensional deployment for high-throughput optical sorters.
\end{abstract}

\section{Introduction}
Commercial dry bean certification and market pricing depend on botanical purity. Hand-sorting grains is laborious and prone to subjective error, necessitating computer vision inspection systems. However, dry bean classification exhibits severe challenges: high dimensional collinearity among geometric descriptors, non-linear boundaries between closely related cultivars (e.g., Dermason vs. Sira), and natural agricultural class imbalance. In this study, we formulate and execute an end-to-end multi-class pipeline conforming to the 11-step ML framework presented in Module~02.

\section{Problem Statement}
We formulate seed cultivar identification as a multi-class inductive classification problem:
\begin{equation}
f: \mathbb{R}^{16} \to \{C_1, C_2, \dots, C_7\}
\end{equation}
where an input vector $\mathbf{x} \in \mathbb{R}^{16}$ contains continuous geometric descriptors, and $y \in \{\text{Dermason, Sira, Seker, Horoz, Cali, Barbunya, Bombay}\}$ is the variety label. Because class frequencies range from 26\% (Dermason) down to 4\% (Bombay), overall classification accuracy alone is misleading. Thus, the primary evaluation criterion is Macro-averaged F1:
\begin{equation}
\text{Macro-F1} = \frac{1}{K} \sum_{k=1}^K \frac{2 \cdot P_k \cdot R_k}{P_k + R_k}
\end{equation}
alongside top-1 accuracy and confusion matrix diagnosis.

\section{Dataset Used}
We utilize the Dry Bean Dataset (UCI Machine Learning Repository ID: 602) \cite{koklu2020multiclass}. The dataset contains 3,500 grain samples characterized by 12 geometric measurements (Area, Perimeter, Major/Minor Axis Length, Convex Area, Equivalent Diameter, Extent, Solidity, Roundness, Compactness) and 4 dimension-free shape factors ($SF_1$ through $SF_4$).

\begin{table}[htbp]
\centering
\caption{Dry Bean Dataset Specification and Partitions}
\label{tab:bean_stats}
\begin{tabular}{ll}
\toprule
\textbf{Specification} & \textbf{Value} \\
\midrule
Total Grains ($N$) & 3,500 instances \\
Target Classes ($K$) & 7 registered bean cultivars \\
Features ($d$) & 16 continuous morphological measurements \\
Class Imbalance & Dermason (26\%), Sira (19\%), Seker (15\%), \\
                & Horoz (14\%), Cali (12\%), Barbunya (10\%), \\
                & Bombay (4\%) \\
Partition Protocol & Stratified 70\% Train ($2,450$) / 15\% Val ($525$) \\
                   & / 15\% Test ($525$ grains) \\
\bottomrule
\end{tabular}
\end{table}

\begin{figure}[htbp]
\centering
\includegraphics[width=\columnwidth]{figures/eda_distribution_and_pca.png}
\caption{Task A3 EDA: (Left) Grain variety class distribution illustrating agricultural skew; (Right) 2D PCA projection demonstrating distinct clusters for Bombay and Seker alongside overlapping Dermason-Sira clusters.}
\label{fig:eda}
\end{figure}

\section{Data Preprocessing and Splitting}
In strict compliance with Module~02's disjoint split mandate, standard scaling parameters ($\boldsymbol{\mu}, \boldsymbol{\sigma}$) are fitted exclusively on the training partition ($N=2,450$) and applied without leakage to validation and test partitions:
\begin{equation}
\tilde{\mathbf{x}} = \frac{\mathbf{x} - \hat{\boldsymbol{\mu}}_{\text{train}}}{\hat{\boldsymbol{\sigma}}_{\text{train}}}
\end{equation}
Stratified sampling preserves cultivar ratios identically across partitions.

\section{Model Architecture and Experimental Setup}
We benchmark five distinct inductive hypotheses:
\begin{enumerate}
    \item \textbf{Multinomial Logistic Regression:} Softmax linear baseline with $\ell_2$ regularization.
    \item \textbf{Balanced Logistic Regression:} Weighting class penalties inversely proportional to training frequency: $w_k = \frac{N}{K \cdot N_k}$.
    \item \textbf{Support Vector Machine (SVM RBF):} Non-linear kernel boundary $K(\mathbf{x}, \mathbf{x}') = \exp(-\gamma \|\mathbf{x} - \mathbf{x}'\|^2)$ with penalty $C = 10.0$.
    \item \textbf{Random Forest Classifier:} Ensemble of 100 decorrelated decision trees with max depth of 12.
    \item \textbf{Multi-Layer Perceptron (MLP):} Fully-connected neural network with hidden layers $(64, 32)$ and ReLU activations.
\end{enumerate}

\section{Results and Performance Analysis}
Test set evaluation on 525 unseen grains is summarized in Table~\ref{tab:results}.

\begin{table}[htbp]
\centering
\caption{Quantitative Performance on Unseen Test Partition ($N=525$)}
\label{tab:results}
\begin{tabular}{lccc}
\toprule
\textbf{Model Architecture} & \textbf{Accuracy} & \textbf{Macro-F1} & \textbf{Weighted-F1} \\
\midrule
Multinomial Logistic Reg. & 85.33\% & 0.8644 & 0.8522 \\
Logistic Reg. (Balanced)  & 84.76\% & \textbf{0.8652} & 0.8479 \\
\textbf{SVM (RBF Kernel)} & \textbf{85.52\%} & 0.8616 & \textbf{0.8541} \\
Random Forest (100 Trees) & \textbf{85.52\%} & 0.8641 & 0.8540 \\
MLP Neural Network        & 84.00\% & 0.8483 & 0.8391 \\
\bottomrule
\end{tabular}
\end{table}

\begin{figure}[htbp]
\centering
\includegraphics[width=\columnwidth]{figures/cm_and_pca_dimension_curve.png}
\caption{Task A3 Diagnostics: (Left) 7x7 confusion matrix for the optimal SVM RBF classifier; (Right) Dimensionality retention curve demonstrating Macro-F1 across PCA components ($k \in \{2, 4, 8, 16\}$).}
\label{fig:cm_pca}
\end{figure}

\section{Ablation and Dimensionality Study}
To evaluate hardware constraints in automated sorting channels, we executed a PCA dimensionality study over components $k \in \{2, 4, 8, 16\}$ using the SVM model. Table~\ref{tab:pca} reports performance retention.

\begin{table}[htbp]
\centering
\caption{PCA Subspace Dimensionality Study (SVM RBF)}
\label{tab:pca}
\begin{tabular}{cccc}
\toprule
\textbf{Components ($k$)} & \textbf{Variance (\%)} & \textbf{Test Macro-F1} & \textbf{\% Retention} \\
\midrule
$k=2$  & 78.4\% & 0.8446 & 98.02\% \\
$k=4$  & 89.1\% & 0.8266 & 95.94\% \\
$k=8$  & 97.2\% & \textbf{0.8543} & \textbf{99.15\%} \\
$k=16$ (Full) & 100.0\% & 0.8616 & 100.00\% \\
\bottomrule
\end{tabular}
\end{table}

Remarkably, retaining only 8 principal axes ($50\%$ feature compression) preserves $99.15\%$ of full classification capability. Even a 2D projection ($k=2$) maintains 0.8446 Macro-F1, illustrating that gross grain scale ($PC_1$) and elongation ($PC_2$) capture the bulk of varietal separability.

\section{Discussion and Error Analysis}
Inspection of the confusion matrix (Fig.~\ref{fig:cm_pca}) confirms domain expectations:
\begin{enumerate}
    \item \textbf{Dermason vs. Sira Overlap:} 30 Sira instances were predicted as Dermason, and 16 Dermason instances were predicted as Sira. In Turkish agronomy, these varieties share nearly indistinguishable circularity ($0.86 \pm 0.03$) and area envelopes.
    \item \textbf{Minority Variety Performance:} Bombay beans (only 21 instances in test) achieved 100\% classification accuracy ($F_1 = 1.000$) due to their gigantic planar area ($>170,000\,\text{pixels}$).
    \item \textbf{Balanced Weighting Efficacy:} Introducing inverse frequency weights improved minority recall on Barbunya from 68\% to 73\%, raising overall Macro-F1 to 0.8652.
\end{enumerate}

\section{Conclusion}
This study verified that non-linear kernel SVMs and Balanced Logistic Regression deliver high-precision dry bean recognition (85.52\% accuracy, 0.8652 Macro-F1). PCA dimensionality reduction showed that an 8-variable representation suffices for real-time sorter deployment with minimal accuracy penalty.

\begin{thebibliography}{1}
\bibitem{koklu2020multiclass}
M.~Koklu and I.~A.~Ozkan, ``Multiclass classification of dry beans using computer vision and machine learning techniques,'' \emph{Computers and Electronics in Agriculture}, vol.~174, p.~105507, 2020.
\bibitem{mitchell1997machine}
T.~M.~Mitchell, \emph{Machine Learning}, McGraw-Hill, New York, 1997.
\bibitem{awais2026introai}
M.~Awais, \emph{AI-101: Introduction to Artificial Intelligence (Course Slides: Modules 01 \& 02)}, Dept. of Intelligent Systems, The University of Lahore, 2026.
\end{thebibliography}

\end{document}
